\begin{table*}[t]\centering\caption{Main results (mean $\pm$ std over 10 seeds, 200 rollouts per run). \texttt{upper\_frac} and \texttt{mode\_balance} are only meaningful targets on \texttt{bimodal80} (0.8) and \texttt{bimodal50} (balance high), respectively; bold/underline marks in the table follow a generic ``higher is better'' convention and should be ignored for these two columns.}\label{tab:main}
\begin{tabular}{llccccc}
\toprule
Method & Task & success $\uparrow$ & collision $\downarrow$ & upper\_frac $\uparrow$ & mode\_balance $\uparrow$ & $n$ \\
\midrule
MSE-MLP & bimodal50 & 0.75 $\pm$ 0.06 & 0.25 $\pm$ 0.06 & \underline{0.52 $\pm$ 0.11} & 0.81 $\pm$ 0.12 & 10 \\
GMM head & bimodal50 & \textbf{1.00 $\pm$ 0.00} & \textbf{0.00 $\pm$ 0.00} & \textbf{0.52 $\pm$ 0.04} & \underline{0.93 $\pm$ 0.05} & 10 \\
\textbf{DDPM head} & bimodal50 & \underline{1.00 $\pm$ 0.00} & \underline{0.00 $\pm$ 0.00} & 0.51 $\pm$ 0.03 & \textbf{0.95 $\pm$ 0.03} & 10 \\
\midrule
MSE-MLP & bimodal80 & 1.00 $\pm$ 0.01 & 0.01 $\pm$ 0.01 & \textbf{1.00 $\pm$ 0.00} & 0.00 $\pm$ 0.00 & 10 \\
GMM head & bimodal80 & \textbf{1.00 $\pm$ 0.00} & \textbf{0.00 $\pm$ 0.00} & 0.81 $\pm$ 0.02 & \textbf{0.38 $\pm$ 0.05} & 10 \\
\textbf{DDPM head} & bimodal80 & \underline{1.00 $\pm$ 0.00} & \underline{0.00 $\pm$ 0.00} & \underline{0.82 $\pm$ 0.02} & \underline{0.36 $\pm$ 0.03} & 10 \\
\midrule
MSE-MLP & unimodal & \textbf{1.00 $\pm$ 0.00} & \textbf{0.00 $\pm$ 0.00} & \textbf{1.00 $\pm$ 0.00} & \textbf{0.00 $\pm$ 0.00} & 10 \\
GMM head & unimodal & \underline{1.00 $\pm$ 0.00} & \underline{0.00 $\pm$ 0.00} & \underline{1.00 $\pm$ 0.00} & \underline{0.00 $\pm$ 0.00} & 10 \\
\textbf{DDPM head} & unimodal & 1.00 $\pm$ 0.00 & 0.00 $\pm$ 0.00 & 1.00 $\pm$ 0.00 & 0.00 $\pm$ 0.00 & 10 \\
\bottomrule
\end{tabular}
\end{table*}

\begin{figure}[t]\centering\includegraphics[width=\columnwidth]{mode_share.pdf}\caption{Upper-mode share of successful rollouts per seed (points) on \texttt{bimodal50} and \texttt{bimodal80}; dashed: demonstrated share.}\label{fig:mode}\end{figure}

\textbf{H1 (MSE fails on \texttt{bimodal50}): supported, in this regime.} MSE-MLP succeeds in \rhval{main/mse-mlp/bimodal50/success/mean:pct1}\% of rollouts versus \rhval{main/gmm-head/bimodal50/success/mean:pct1}\% (GMM) and \rhval{main/ddpm-head/bimodal50/success/mean:pct1}\% (DDPM); the DDPM--MSE difference is \rhval{cmp/main/mse-mlp/bimodal50/success/delta:3} (Welch $p=\rhval{cmp/main/mse-mlp/bimodal50/success/welch_p}$, Cohen's $d=\rhval{cmp/main/mse-mlp/bimodal50/success/cohen_d:1}$). Essentially all MSE failures are collisions (collision rate \rhval{main/mse-mlp/bimodal50/collision/mean:pct1}\%), consistent with trajectories that head partly toward the obstacle. The size of the effect depends strongly on start jitter (Section~\ref{sec:ablations}), so ``MSE fails'' is a statement about near-identical starts, not about regression in general.

\textbf{H2 (DDPM does not beat GMM): supported (no difference found).} The DDPM--GMM success difference is \rhval{cmp/main/gmm-head/bimodal50/success/delta:3} (Welch $p=\rhval{cmp/main/gmm-head/bimodal50/success/welch_p}$); both are at ceiling, so success cannot discriminate between them here. Mode balance on \texttt{bimodal50} is also similar (difference \rhval{cmp/main/gmm-head/bimodal50/mode_balance/delta:3}, $p=\rhval{cmp/main/gmm-head/bimodal50/mode_balance/welch_p}$). A ceiling effect means this task cannot reveal an advantage of DDPM; it is not evidence that none exists on harder tasks.

\textbf{H3 (mode-share fidelity on \texttt{bimodal80}): supported for GMM and DDPM; MSE collapses to one mode.} Upper share among successes is \rhval{main/gmm-head/bimodal80/upper_frac/mean:2} (GMM) and \rhval{main/ddpm-head/bimodal80/upper_frac/mean:2} (DDPM) against a demonstrated 0.8, both within the pre-registered 0.1 (difference $p=\rhval{cmp/main/gmm-head/bimodal80/upper_frac/welch_p}$). MSE-MLP has upper\_frac \rhval{main/mse-mlp/bimodal80/upper_frac/mean:2}: it always passes on the majority side and still succeeds in \rhval{main/mse-mlp/bimodal80/success/mean:pct1}\% of rollouts, because with a 0.8 split the average waypoint ($y\approx0.21$) clears the obstacle. Thus MSE is a good \emph{task-success} policy on \texttt{bimodal80} but does not reproduce the demonstrated diversity; whether diversity matters depends on the application. On the \texttt{unimodal} control all three heads reach \rhval{main/mse-mlp/unimodal/success/mean:pct1}\%, \rhval{main/gmm-head/unimodal/success/mean:pct1}\% and \rhval{main/ddpm-head/unimodal/success/mean:pct1}\%, so the failures above are due to multimodality and not to a capacity or training problem in the setup.

\textbf{Where the methods fail.} Beyond MSE on \texttt{bimodal50}, the DDPM head's failures are in the ablations: one denoising step and long action chunks (Section~\ref{sec:ablations}). The study cannot say anything about higher-dimensional actions, image inputs or multi-step stochastic sampling effects, since none of the metrics here separate the two multimodal heads.
